% ---------------------------------------------------------------------------
% Step 1 of the physics chain: the exact hierarchy.
%
% Written 2026-10-24 by the Paper 1 stream. The supplement had no section for
% the object every other section is about -- a reader met tau_{A_l} and
% (a B_1)' with no account of the hierarchy they belong to.
%
% Equations read from the antecedent and from functions/spectra/. Findings
% cited, none minted.
% ---------------------------------------------------------------------------
\section{The exact hierarchy}
\label{sec:hierarchy}

\emph{This is the object every later section is about. It is written once,
here, with what ``exact'' is claiming and what it is not.}

% ===========================================================================
\subsection{The expansion, and why it is PSTF}

The photon distribution is expanded in \textbf{projected, symmetric,
trace-free} temperature multipoles about a chosen 4-velocity $u^a$:
\begin{equation}
  \tau(x,e)=\sum_{\ell\ge0}\tau_{A_\ell}(x)\,O^{A_\ell}(e),
  \qquad
  O^{A_\ell}=e^{\langle A_\ell\rangle},
  \label{eq:pstf-expansion}
\end{equation}
with $A_\ell$ the multi-index $a_1\ldots a_\ell$, $e^a$ a unit direction
orthogonal to $u^a$, and $\langle\,\rangle$ the PSTF part.

\textbf{PSTF is not a choice of basis among several.} On the 2-sphere orthogonal to
$u^a$ the PSTF tensors are the \emph{irreducible} representations of the
rotation group: each rank carries exactly the information no lower rank
carries, and the decomposition is unique. Any other basis --- spherical
harmonics, Legendre polynomials in $\hat k\cdot e$ --- is a representation of
the same content that first requires a direction to be singled out.
\textbf{The PSTF form singles out nothing but $u^a$}, which is the whole
reason a $1+3$ treatment can carry the frame generically and choose it at the
end.

\paragraph{Two objects, and conflating them is how a sign is lost.}
$\tau_{A_\ell}$ are \textbf{multipole} coefficients and exist for a general
geometry. $\tau_\ell$ are \textbf{mode} coefficients, obtained by expanding
$\tau_{A_\ell}$ on harmonics $Q_{A_\ell}$, and exist only for an
almost-Robertson--Walker one. \emph{Annals~II} notes the distinction is not
made in Bardeen-variable treatments. \textbf{Every released figure plots the
multipole quantity}, and criterion~3 --- the covariant-to-mode round trip ---
is where the two are held to each other.

% ===========================================================================
\subsection{The transport equation}

With $\kappa'=an_e\sigma_{\rm T}$ the differential optical depth, the exact
covariant transport equation at rank $\ell$ is \emph{Annals~II}~(68):
\begin{equation}
  \dot\tau_{A_\ell}
  +\underbrace{\tfrac{\ell+1}{2\ell+3}\mathrm D^a\tau_{A_\ell a}
               -\mathrm D_{\langle a_\ell}\tau_{A_{\ell-1}\rangle}}_{\text{free streaming}}
  +\underbrace{\kappa'\tau_{A_\ell}}_{\text{out-scattering}}
  =\underbrace{a\mathcal B_{A_\ell}}_{\text{gravitational}}
  +\underbrace{\kappa'\mathcal S_{A_\ell}}_{\text{in-scattering}} .
  \label{eq:transport-exact}
\end{equation}

\textbf{Four terms, and each later section is about one of them.} Free
streaming is \S\ref{sec:free-streaming}; the gravitational source closes at
rank two and is \S\emph{The generic source derivatives}; the two $\kappa'$
terms are the two halves of one collision term and their separation is what
\S\ref{sec:w-form} turns on.

\paragraph{Both $\kappa'$ are the same $\kappa'$.} The out-scattering on the
left removes photons from the beam and the in-scattering on the right returns
them; they are one Thomson process written on two sides. \textbf{That is why
multiplying through by the integrating factor assigns both weights at once
rather than by choice}, which is \eqref{eq:w-form}'s content.

% ===========================================================================
\subsection{What ``exact'' claims, and what it does not}

\textbf{Claimed.} \eqref{eq:transport-exact} is exact in the $1+3$ split: no
expansion in a perturbation parameter has been made in deriving it, no
background has been subtracted, and no gauge has been fixed. It holds for a
general $u^a$ in a general geometry.

\textbf{Not claimed.} That the \emph{solution} computed here is exact. Every
approximation this release makes enters after \eqref{eq:transport-exact} and
is named --- the scattering-time expansion, \textbf{A-3}'s static transfer,
\textbf{A-5}'s collision term, the diffusion treatment --- and
\S\ref{sec:open-questions} carries what is not costed.

\paragraph{The sources are gauge-invariant, and for a reason worth stating
once.} $\mathcal B_0$, $\mathcal B_a$ and $\mathcal B_{ab}$ all vanish in the
background, so by \textbf{Stewart--Walker} they are gauge-invariant at first
order --- which is the paper's own \emph{``note that these are CGI with
respect to $u^a$''}. \textbf{Gauge-invariance here is a consequence of the
objects chosen, not a construction imposed on them}, and that is the
difference between a covariant treatment and a gauge-fixed one with covariant
notation.

\paragraph{What the collision term contains, and what A-5 drops.}
\emph{Annals~II} writes
$C[\tau]\simeq\sigma_{\rm T}n_E\big(\bar f_1e^av^B_a+\bar f_2e^ae^b\tau_{ab}
-\tau\big)$. This release takes $\bar f_1=1$ and $\bar f_2=0$, so
$\mathcal S_2=0$: \textbf{no polarisation feedback and no quadrupole
in-scattering}. That is approximation \textbf{A-5} stated at the collision
term rather than at the oscillator, and \S\ref{sec:w-form} shows it cannot
produce a frame disagreement at linear order.

% ===========================================================================
\subsection{Where the hierarchy is gated}

\begin{center}\small
\begin{tabularx}{\textwidth}{@{}l >{\raggedright\arraybackslash}X l@{}}
\toprule
\textbf{gate} & \textbf{what it tests} & \textbf{what it does not} \\
\midrule
\textbf{G1}, criterion~1 & the free-streaming part of
  \eqref{eq:transport-exact} against four external hierarchies, two
  normalisations, two phase conventions & anything with the source on ---
  \S\ref{sec:free-streaming} \\
\texttt{test\_register} & that every module, object and test named in the
  equation and notation tables resolves & that the equation implemented is
  the one printed \\
\textbf{C2} round trip & covariant $\to$ mode $\to$ covariant, identically
  zero in exact arithmetic & that either representation is the right one ---
  that is criterion~3, reopened \\
\bottomrule
\end{tabularx}
\end{center}

\noindent
\textbf{No gate tests \eqref{eq:transport-exact} as a whole}, and none can:
it is the antecedent's equation, read from the page. Its standing is
\emph{inspection} --- \S\emph{the acceptance spine} ---  and every numerical
claim in this supplement is about what was done \emph{with} it.
